\bookchapter{פרק 52: ארבעת הרכיבים ומחזור החיים של Activity}

\chapterintro{אם Compose הוא “הפנים” של האפליקציה, הפרק הזה עוסק ב“שלד” שלה. למערכת אנדרואיד יש ארבעה רכיבים בסיסיים, וכל אפליקציה בנויה מהם; והחשוב מביניהם, Activity, מלווה ב“מפת מחזור חיים” שכל מפתח אנדרואיד חייב לחרוט בזיכרון. אחרי הפרק הזה תבינו חידות נצחיות כמו “למה כשחוזרים מהרקע, המצב של האפליקציה שלי נעלם”.}

\booksection{א' ארבעת הרכיבים: עמודי התווך של אנדרואיד}

אנדרואיד מגדירה רשמית ארבעה רכיבים בסיסיים, שכל אחד מהם אחראי על סוג עבודה אחר:

\begin{center}
\begin{tabularx}{\linewidth}{|X|X|X|}
\hline
\textbf{רכיב} & \textbf{תפקיד} & \textbf{אנלוגיה} \tabularnewline
\hline
Activity & מסך גלוי לעין & חזית החנות – המשתמש רואה אותה ומקיים איתה אינטראקציה \tabularnewline
Service & "הבלתי נראה" שעובד ברקע & פועל המחסן – מנגן מוזיקה, מוריד קבצים; לא צריך מסך \tabularnewline
BroadcastReceiver & קבלת שידורי מערכת & חדר השער – הרמקול שמכריז "הסוללה נגמרה" ו"הרשת התנתקה" \tabularnewline
ContentProvider & שיתוף נתונים בין אפליקציות & ארון מסמכים ציבורי – למשל לשתף אנשי קשר עם אפליקציות אחרות \tabularnewline
\hline
\end{tabularx}
\end{center}

\begin{center}\small\color{gray}סקירה מהירה של ארבעת רכיבי אנדרואיד\end{center}

בשלב המתחיל תבלו כ-95\% מהזמן עם Activity; ‏Service ו-BroadcastReceiver נדרשים לנגינת מוזיקה ולדחיפת התראות; את ContentProvider תשתמשו הכי פחות (המערכת כבר עוטפת את הנתונים הנפוצים). שימו לב: כל אחד מארבעת הרכיבים חייב “להירשם” ב-AndroidManifest.xml – בלי רישום, המערכת פשוט לא תמצא אותם.

\booksection{ב' מחזור החיים של Activity: “מפת האוצר” הראשונה בפיתוח אנדרואיד}

אפליקציה בטלפון עלולה להיקטע בכל רגע: שיחה נכנסת, המשתמש עובר לענות ב-WeChat, למערכת אין מספיק זיכרון והיא מסיימת את התהליך שלך... חיי Activity הם מסע מעגלי על המפה הזאת:

onCreate (לידה: אתחול)

onStart (גלוי)

onResume (בחזית, אינטראקציה)

onPause (מוסתר חלקית)

onStop (בלתי נראה לגמרי)

onDestroy (הריסה)

onRestart / onStart / onResume: בחזרה לחזית, אותו מסלול בכיוון הפוך

\begin{center}\small\color{gray}איור 4: מחזור החיים של Activity (כיוון ההתקדמות: השורה העליונה – לידה השורה התחתונה – יציאה; בדרך חזרה – אותו נתיב בכיוון ההפוך)\end{center}

\begin{center}
\begin{tabularx}{\linewidth}{|X|X|X|}
\hline
\textbf{פונקציית Callback} & \textbf{מתי מופעלת} & \textbf{שימוש טיפוסי} \tabularnewline
\hline
onCreate & בעת יצירת ה-Activity (פעם אחת) & טעינת המסך ואתחול הנתונים \tabularnewline
onStart & המסך הופך לגלוי & התחלת לוגיקה קלה \tabularnewline
onResume & מקבל מיקוד ומוכן לאינטראקציה & התחלת אנימציות, המשך נגינה \tabularnewline
onPause & ברגע הכיסוי החלקי או המעבר למסך אחר & השהיית אנימציות, שמירת טיוטה \tabularnewline
onStop & בלתי נראה לגמרי & שחרור משאבים כבדים, השהיית הורדות \tabularnewline
onDestroy & לפני ההריסה & ניקוי אחרון \tabularnewline
\hline
\end{tabularx}
\end{center}

\begin{center}\small\color{gray}השימושים הטיפוסיים של שש פונקציות ה-Callback של מחזור החיים\end{center}

\textbf{אזהרה: סיבוב המסך “הורג” את ה-Activity}\par

המקרה הקלאסי של מתחילים: מסובבים את המסך וכל הטקסט בשדות הקלט נעלם. הסיבה: סיבוב מסך כברירת מחדל הורס ובונה מחדש את ה-Activity (עובר על onDestroy onCreate). הפתרון: שמירת מצב עם rememberSaveable (התחליף ל-remember ב-Compose), שמשחזר את הערכים אוטומטית בעת הבנייה מחדש. נשתמש בזה שוב ושוב בפרקים 6 ו-7.

\booksection{ג' Intent: “כרטיס הנסיעה” לתנועה בתוך האפליקציה ומחוצה לה}

‏Intent (כוונה) הוא “כרטיס הנסיעה” לתקשורת בין רכיבים. כדי לקפוץ ממסך אחד לאחר, לפתוח דפדפן או להפעיל מצלמה – צריך “לקנות כרטיס ולעלות לרכבת”:

\begin{lstlisting}[language=Kotlin]
// דרך 1: Intent מפורש – מציינים בבירור לאיזה מסך ללכת
val intent = Intent(this, DetailActivity::class.java)
intent.putExtra("name", "דני")     // לוקח "מזוודות" לדרך (הנתונים)
startActivity(intent)

// דרך 2: Intent גלוי-מרומז – אומרים רק מה לעשות והמערכת מוצאת מי יעשה
val web = Intent(Intent.ACTION_VIEW, Uri.parse("https://developer.android.com"))
startActivity(web)   // המערכת תפתח את הדפדפן (או תבקש ממך לבחור)
\end{lstlisting}

לשתי הדרכים שימושים שונים: מפורשת – למעברים בין מסכים בתוך האפליקציה שלכם; מרומזת – להפעלת יכולות של המערכת או של אפליקציות אחרות (פתיחת דף אינטרנט, שיחה, שיתוף). בפרויקט Compose, בתוך Composable משיגים את ה-Context עם LocalContext.current, ואז אפשר להריץ startActivity.

\booksection{ד' מספר מסכים וניווט}

לכל אפליקציה אמיתית יש כמה מסכים. פתרון הניווט הרשמי המומלץ ל-Compose הוא Navigation Compose (ספריית התלות androidx.navigation):

\begin{lstlisting}[language=Kotlin]
val navController = rememberNavController()

NavHost(navController, startDestination = "home") {
    composable("home") { HomeScreen(navController) }
    composable("detail/{name}") { entry ->
        val name = entry.arguments?.getString("name")
        DetailScreen(name ?: "לא ידוע")
    }
}

// כדי לנווט מתוך המסך:
navController.navigate("detail/דני")     // עם פרמטר
navController.popBackStack()              // חזרה למסך הקודם
\end{lstlisting}

\begin{tcbinfo}{טיפ: מה עושים כשהמסכים מתרבים?}
כשהפרויקט גדל, רכזו את כל נתיבי המסכים בקובץ אחד (למשל AppNavGraph.kt) לניהול אחיד – אל תפזרו את הלוגיקה של “לאן עוברים” בין המסכים. זהו הצעד הראשון מ“רץ” אל “ניתן לתחזוקה”.

\end{tcbinfo}

\booksection{ה' בדרך אגב: מהו Fragment (ולמה הספר הזה לא מלמד אותו)}

בחיפוש אחר מדריכים ישנים בוודאי תיתקלו ב-Fragment – הוא “קטע מסך לשימוש חוזר” בתוך Activity, והיה הכוכב הראשי בארכיטקטורות רבי-מסכים ישנות. אבל בעידן ה-Compose מספיק לבנות מסכים ישירות עם פונקציות Composable + נתיבי Navigation, וגם הרשויות הפחיתו בחשיבותו של Fragment. לכן הספר הזה לא כולל אותו בקו הראשי.

ההמלצה שלנו: בפרויקט חדש אל תשתמשו ב-Fragment מיוזמתכם; אבל בקריאת קוד ישן או בראיון עבודה, מספיק שתדעו שזהו “קטע מסך משובץ בתוך Activity”. במשאבי לימוד תעדיפו תמיד מ-2023 ואילך – “לוח השנה הישן” של טכנולוגיית האנדרואיד מתעדכן מהר יותר משחושבים.

\begin{tcbpractice}{הגיע הזמן להתנסות: תנו למונה "זיכרון"}
שדרגו את המונה מפרק 51: ‏(1) החליפו את remember ב-rememberSaveable וסובבו את המסך כדי לוודא שהמספר נשאר; (2) הוסיפו כפתור “אודות” שמנווט עם Navigation למסך שני שמציג את מספר הגרסה; (3) הדפיסו לוגים (Log.d) ב-onCreate/onPause של שני המסכים וצפו בסדר הקריאות של מחזור החיים דרך Logcat. אחרי שלושת הצעדים האלה, מחזור החיים כבר לא יישאר תיאוריה על הנייר.

\end{tcbpractice}

\begin{tcbsummary}{סיכום הפרק}
\begin{itemize}
  \item ארבעת הרכיבים: Activity (חזית), Service (רקע), BroadcastReceiver (רמקול), ContentProvider (ארון נתונים)
  \item שש פונקציות ה-Callback של מחזור החיים: onCreate onStart onResume onPause onStop onDestroy
  \item סיבוב מסך הורס ובונה מחדש את ה-Activity: שמרו מצב עם rememberSaveable
  \item ‏Intent הוא "כרטיס הנסיעה": מפורש – למסכים שלכם; מרומז – ליכולות המערכת
  \item ריבוי מסכים עם Navigation Compose: NavHost + נתיבי composable
\end{itemize}

\end{tcbsummary}
